\section{A Measurement Model for Benchmark Reproduction}
\label{sec:model}

\subsection{The generative model of a score}

Let a measured benchmark result be
\begin{equation}
Y \;=\; F(R,\,D,\,I,\,A,\,U),
\label{eq:model}
\end{equation}
where
\begin{itemize}[nosep,leftmargin=1.4em]
  \item $R$ --- the \textbf{recorded recipe}: scenario, dataset, prompt
        construction, in-context examples, decoding parameters, and scorer;
  \item $D$ --- the \textbf{deployment}: weights, model and tokenizer revisions,
        chat template, load precision (dtype), and quantization;
  \item $I$ --- the \textbf{execution instrument}: harness version, serving
        engine, package and kernel versions, hardware, device topology, and
        batching/caching;
  \item $A$ --- the \textbf{artifact interpretation}: storage schema, migration,
        canonicalization, and carried-forward outputs; and
  \item $U$ --- \textbf{residual} stochastic and numerical influences.
\end{itemize}
The public \helm record exposes $R$ and the \emph{names} within $D$, but only a
projection of $D$'s values and almost nothing of $I$, $A$, or $U$. Reproduction
is thus an inverse problem: infer whether a locally chosen $(\hat D, \hat I)$
yields $F(R,\hat D,\hat I,\hat A,U)\approx Y$, and if not, \emph{which coordinate}
is responsible. We say a run is \emph{identifiable at tolerance $\epsilon$} when
all materially compatible configurations produce equivalent results under a
declared output-, metric-, or conclusion-level criterion.

\paragraph{The taxonomy this model implies.} Equation~\ref{eq:model} does not
merely describe a score; it dictates how reproduction attempts can be
classified, and that taxonomy organizes the rest of the paper. Whether a run can
be executed at all is a precondition rather than an outcome, and forms an
upstream gate (\S\ref{sec:gate}). Past that gate, an attempt is characterized
along two orthogonal axes: which coordinates of the model it puts under test
(\S\ref{sec:targets}), and, when its result disagrees with the record, which
coordinate is responsible (\S\ref{sec:taxonomy}). Each axis is read off the
model rather than proposed independently of it --- the first partitions the
coordinates into those a given attempt exercises and those it holds fixed, the
second assigns a disagreement to one coordinate, a residual, or an
underdetermined case. Neither axis refines the other, and collapsing them is
what produces the familiar but uninformative verdict that a benchmark ``did not
reproduce.''

\subsection{The eligibility gate}
\label{sec:gate}

Before any disagreement is measured, a run must be executable. Gated models or
datasets, revoked download access, missing credentials, unsupported packaging,
unavailable judges, and infeasible resource requirements are
\emph{eligibility/environment filters}. They answer ``can we run this?'', not
``did it reproduce?'', and must never be folded into a negative reproduction
result. Making this gate explicit --- every excluded run carries a typed
exclusion reason --- is what allows a defensible ``$X$ of $Y$ runnable runs''
denominator.

The gate is relative to a target, not to a run. The sharpest case is a
model-graded scenario whose judge is no longer served: the judging step is
unrepeatable, so the run is ineligible for procedural reproduction, but the
judge's per-instance verdicts are themselves part of the released record, and
the published aggregate can still be recomputed from them and checked. Such a
run is out of scope for the second target of \S\ref{sec:targets} and fully in
scope for the first. Treating eligibility as a property of the run alone would
discard precisely the evidence that survived, and would understate the auditable
fraction of the corpus; we therefore scope every exclusion reason to the target
it blocks.

\subsection{Axis 1: what is under test}
\label{sec:targets}

A reproduction attempt must first declare what it is attempting to reproduce.
We adapt the standard methods/results/inferential trichotomy of
\citet{goodman2016does} to a benchmark corpus that publishes
per-instance outputs. The three resulting choices each exercise a different
subset of the coordinates of Equation~\ref{eq:model}, and they are not
interchangeable:
\begin{enumerate}[nosep,leftmargin=1.6em]
  \item \textbf{Artifact reconstruction}: recompute the published aggregates from
        the released prompts, completions, and metric records \emph{without}
        rerunning a model. Tests $A$ and the scorer only.
  \item \textbf{Procedural reproduction}: execute the recorded recipe under a
        pinned or explicitly equivalent instrument and compare instance-level
        outputs and metrics. Tests $D$ and $I$.
  \item \textbf{Claim-level robustness}: determine whether the substantive
        conclusion (a ranking, or a model-generation improvement) survives
        plausible execution substrates even when byte-identical outputs do not.
\end{enumerate}
Targets (2) and (3) are Goodman's methods and inferential reproducibility --- the
former with its ``same data and tools'' premise replaced by a declared-equivalent
instrument, since the original tools are what did not survive. Target (1) has no
counterpart there: it is available only because \helm releases the per-instance
prompts, completions, and metric records that make a score recomputable without
executing a model. Nor is it a fallback for runs that are merely inconvenient to
execute --- for model-graded scenarios whose judge is no longer served it is the
only target available at all, and it still answers a real question about them,
since a scorer or migration defect would show up here with no judge in the loop. Goodman's results reproducibility, which asks for corroboration
from an independent study, has no counterpart here in the other direction --- a
fixed historical corpus admits no new study.

Released official artifacts remain the historical ground truth for what \helm
reported; a rerun assesses the \emph{stability and identifiability} of the
procedure that produced them, and does not overwrite the record. This distinction
is what licenses us to report a differing rerun as evidence about $I$ rather than
as a correction of the official number.

\subsection{Axis 2: why two executions disagree}
\label{sec:taxonomy}

Once the target is fixed and the attempt executes, an observed disagreement
admits one of six causes --- one per coordinate of Equation~\ref{eq:model}, plus
a residual and an underdetermined case:
\begin{enumerate}[nosep,leftmargin=1.6em]
  \item \textbf{Recipe mismatch} --- $R$ differs (prompt construction, examples,
        decoding, scorer).
  \item \textbf{Deployment mismatch} --- $D$'s client/endpoint/engine identity
        differs (hosted API vs.\ local; in-process \code{transformers} vs.\ vLLM).
  \item \textbf{Execution-instrument mismatch} --- $I$ differs (precision,
        tokenizer/template version, attention kernel, device placement,
        software-stack versions).
  \item \textbf{Artifact-interpretation/migration mismatch} --- $A$ differs
        (row reordering, schema migration, carried-forward outputs).
  \item \textbf{Residual disagreement under controlled equivalence} --- all
        controllable layers are matched and a gap remains. This is the only
        category analogous to conventional repeatability failure.
  \item \textbf{Historically non-identifiable reproduction} --- surviving
        evidence does not uniquely specify what execution should count as
        faithful. Qualitatively different from (5): not ``it failed to repeat''
        but ``what would count as a faithful execution is underdetermined.''
\end{enumerate}
The two axes are independent: any target can fail for any cause, and it is their
combination, not either alone, that describes an outcome. We deliberately avoid
the phrase ``true reproducibility failure'' whenever execution provenance is
incomplete; categories (1)--(4) are usually recoverable by controlling the
substrate, and (6) is a first-class result, not a gap to be closed.

\subsection{Per-parameter identifiability}
\label{sec:identifiability}

Identifiability is not a binary property of an entire run: each latent parameter
of $(D,I)$ has its own status, and the paper's per-case analysis assigns one to
each (Table~\ref{tab:identifiability}). This refinement is what makes both the
positive attributions and the non-identifiability result precise: we never claim
to identify a whole historical deployment, only specific coordinates, each with
its own evidentiary basis.

\begin{table}[t]
\centering\small
\begin{tabularx}{\linewidth}{@{}p{0.26\linewidth}p{0.20\linewidth}Y@{}}
\toprule
\textbf{Parameter} & \textbf{Typical status} & \textbf{Basis / limitation} \\
\midrule
Run specification ($R$) & identifiable & frozen public \file{run\_spec.json};
migration can still alter interpretation. \\
Client / deployment class & often identifiable & recorded deployment name and
harness config; hosted-provider internals remain opaque. \\
Load dtype & conditional & explicit when pinned; else inferable only from loader
source $+$ package-version bounds. \\
Tokenizer special-token behavior & often recoverable & repo/config inspection and
token-level prompt comparison. \\
Chat template / generation prompt & often recoverable & revisioned tokenizer files
or controlled rendering sweeps; effective behavior depends on package version. \\
Model / tokenizer revision & frequently missing & rarely pinned; mutable aliases
prevent unique historical recovery. \\
Engine / package / kernel versions & often missing & client family may be known;
exact versions/kernels usually not. \\
Batching, cache, topology, hardware & usually missing historically & measure
robustness rather than promise exact replay. \\
Harness era & tagged or proxy-only & tagged releases containerizable; migrated
paths may admit only declared proxy instruments. \\
\bottomrule
\end{tabularx}
\caption{Per-parameter identifiability in the studied cases. The status is
\emph{per coordinate}, not per run; \S\ref{sec:cases} assigns each case's
disagreement to a coordinate and its status here.}
\label{tab:identifiability}
\end{table}
